\documentclass[11pt]{article}
\usepackage[margin=1in]{geometry}
\usepackage[T1]{fontenc}
\usepackage{lmodern,microtype,amsmath,amssymb,amsthm,mathtools}
\usepackage[colorlinks=true,linkcolor=blue,citecolor=blue,urlcolor=blue]{hyperref}
\newtheorem{theorem}{Theorem}[section]
\newtheorem{lemma}[theorem]{Lemma}
\newtheorem{proposition}[theorem]{Proposition}
\newtheorem{corollary}[theorem]{Corollary}
\newcommand{\E}{\mathbb E}
\newcommand{\R}{\mathbb R}
\newcommand{\Hh}{\mathcal H}
\newcommand{\HS}{\mathrm{HS}}
\newcommand{\tr}{\operatorname{tr}}
\newcommand{\Cov}{\operatorname{Cov}}
\title{Completed Gaussian quadratic products:\protect\\
sharp fourth moments and contraction tests}
\author{Project Stardust\\\small A companion to a frozen verified library}
\date{October 29, 2026}
\begin{document}
\maketitle
\begin{abstract}
Centered Gaussian quadratic forms extend from finite matrices to the
complete second Wiener chaos. We give complete proofs of the sharp
dimension-independent fourth-moment bounds, bounded multiplication
into Gaussian $L^2$, its exact norm, and the degree-zero/two/four
product decomposition. A nonzero completed quadratic has strictly
larger fourth moment than a Gaussian of the same variance, although
the lower universal coefficient is not attained. Polarization gives
nonnegative covariance of quadratic squares with exact contraction
conditions for equality. For finite symmetric matrices these conditions
reduce to matrix annihilation, yielding the classical independence
test. The completion and its source/operator normalization are made
explicit. Classical quadratic-form theory is credited throughout;
no new independence criterion or fourth-moment phenomenon is claimed.
\end{abstract}

\section{The actual Gaussian space and the main statements}
Let $(G_i)_{i\geq0}$ be independent standard real Gaussian variables
on their product probability space. Every $L^p$ space below uses
this positive probability law. Let $\Hh_2$ be the closed subspace of
$L^2$ generated by the centered quadratic monomials
\[
 G_i^2-1,\qquad G_iG_j\quad(i<j).
\]
This is the second Wiener chaos. Orthogonality of the tensor Hermite
polynomials also defines the closed degree-$k$ spaces $\Hh_k$ and
their orthogonal projections $P_k$. We use the probabilists'
normalization $H_2(x)=x^2-1$, $H_4(x)=x^4-6x^2+3$.
Here $\Hh_0$ is the constant space and $P_0F=(\E F)1$.

\begin{theorem}\label{thm:main}
Every $F\in\Hh_2$ belongs to $L^4$. If $F\ne0$, then
\begin{equation}\label{eq:mainbounds}
 3(\E F^2)^2<\E F^4\leq15(\E F^2)^2.
\end{equation}
The upper coefficient is attained. The lower coefficient $3$ is
optimal over the entire space, but equality there forces $F=0$.
Multiplication is a bounded bilinear map
\[
 (F,G)\longmapsto FG:\Hh_2\times\Hh_2\to L^2,
 \qquad \|FG\|_2\leq\sqrt{15}\|F\|_2\|G\|_2,
\]
whose norm is exactly $\sqrt{15}$. Its range lies in
$\Hh_0\oplus\Hh_2\oplus\Hh_4$, rather than in $\Hh_2$ alone.
The exact fourth-cumulant identity is
\begin{equation}\label{eq:cumulant}
 \E F^4-3(\E F^2)^2=\frac32\|P_2(F^2)\|_2^2.
\end{equation}
For all $F,G\in\Hh_2$,
\begin{align}\label{eq:polarized}
 \Cov(F^2,G^2)
 &=2\langle F,G\rangle^2
 +\frac12\langle P_2(F^2),P_2(G^2)\rangle
 +\|P_2(FG)\|_2^2\\
 &\geq2\langle F,G\rangle^2\geq0.\nonumber
\end{align}
Every term on the first line's right side is nonnegative. In
particular equality $\Cov(F^2,G^2)=0$ holds precisely when all
three displayed terms vanish.
\end{theorem}

The assertion about the middle contraction pairing is special to
quadratic chaos. It is not a generic assertion that inner products
of arbitrary orthogonal projections are nonnegative. Likewise,
uncorrelated quadratic observables need not be independent.

Lancaster's primary quadratic-form work \cite{Lancaster} identifies
cumulants with matrix traces. Ogawa's primary account \cite{Ogawa}
states the classical centered Gaussian matrix-annihilation
independence criterion. The present paper reconstructs the exact
completed-space mechanisms in a verified library; it does not claim
priority for those classical results. Primary full-text access for
Lancaster's detailed hypotheses remains incomplete.

\section{Finite matrices, observable normalization and moments}
For a real symmetric $n\times n$ matrix $A$ define
\begin{equation}\label{eq:finiteQ}
 Q_A=\sum_{i,j<n}A_{ij}(G_iG_j-\delta_{ij})
     =G^T A G-\tr A.
\end{equation}
If a matrix is not symmetric, its observable equals that of
$S(A)=(A+A^T)/2$: transpose the double sum to see that the
antisymmetric part vanishes. Thus the source observable identifies
matrices differing by a skew-symmetric matrix. We first work on the
symmetric representative to keep its norm nondegenerate.

The generating function of independent standard Gaussian coordinates
is $\E e^{\langle u,G\rangle}=e^{\|u\|^2/2}$. Expanding it to
degree four gives
\[
 \E G_iG_jG_kG_l
 =\delta_{ij}\delta_{kl}+\delta_{ik}\delta_{jl}
       +\delta_{il}\delta_{jk}.
\]
Consequently,
\begin{equation}\label{eq:isometry}
 \E Q_A=0,\qquad
 \langle Q_A,Q_B\rangle=2\sum_{i,j}A_{ij}B_{ij},\qquad
 \E Q_A^2=2\|A\|_\HS^2.
\end{equation}
The Hilbert--Schmidt norm here is the Frobenius norm on finite
matrices, $\|A\|_\HS^2=\sum_{i,j}A_{ij}^2$.

An orthogonal change of coordinates preserves the standard Gaussian
law, since it preserves its characteristic function
$e^{-\|u\|^2/2}$. Diagonalize $A$ by a real orthogonal matrix. If
its eigenvalues are $a_i$, then $Q_A$ has the law of
$\sum_i a_i(G_i^2-1)$. Set $U_i=G_i^2-1$. The elementary normal
moments $\E G^{2k}=(2k-1)!!$ follow by Gaussian integration by parts,
and give
\[
 \E U_i=0,\qquad \E U_i^2=2,\qquad \E U_i^4=60.
\]
In the fourth power of the independent sum, only the four-equal
and two-pair index patterns survive. Thus
\begin{align}\label{eq:finitefourth}
 \E Q_A^4
 &=60\sum_i a_i^4+24\sum_{i<j}a_i^2a_j^2\\
 &=12\Bigl(\sum_i a_i^2\Bigr)^2+48\sum_i a_i^4.
 \nonumber
\end{align}
Since $\sum_i a_i^4\leq(\sum_i a_i^2)^2$, this proves the finite
upper bound $\E Q_A^4\leq15(\E Q_A^2)^2$. Also
\begin{equation}\label{eq:finitecumulant}
 \E Q_A^4-3(\E Q_A^2)^2
 =48\|A^2\|_\HS^2,
\end{equation}
because the eigenvalues of $A^2$ are $a_i^2$. No simultaneous
diagonalization of two unrelated symmetric matrices is used.

\section{The full second chaos and its original source field}
Let $\mathcal S_2$ be the real Hilbert space of symmetric square-summable
matrices $A=(A_{ij})_{i,j\geq0}$, with its Hilbert--Schmidt norm.
Let $A^{(n)}$ retain its first $n$ rows and columns and set all other
entries to zero. Then $A^{(n)}\to A$ in that norm. By
\eqref{eq:isometry}, the actual random variables $Q_{A^{(n)}}$
are Cauchy in $L^2$. Define
\[
 Q_A=\lim_{n\to\infty}Q_{A^{(n)}}\quad\hbox{in }L^2.
\]
The resulting map satisfies
\begin{equation}\label{eq:fullisometry}
 \|Q_A\|_2^2=2\|A\|_\HS^2.
\end{equation}
It is onto $\Hh_2$: its range is closed by this isometry and
contains every centered quadratic monomial. Conversely, every
finite-box observable lies in the closed span defining $\Hh_2$.
This proves both inclusions and uniqueness of the symmetric matrix.
The definition does not require $\sum_i A_{ii}$ to converge.
The formal expression $G^TAG-\tr A$ is only the finite formula;
the completed centered quadratic is an $L^2$ limit, not an
unjustified subtraction of two infinite quantities.

The same matrix has an actual Gaussian linear source field. Let
$H=\ell^2(\mathbb N)$ with orthonormal basis $(e_i)$ and define
\[
 V_A=\lim_{n\to\infty}\sum_{i,j<n}A_{ij}e_iG_j
       \quad\hbox{in }L^2(\Omega;H).
\]
Orthonormality gives $\|V_A\|^2=\|A\|_\HS^2$. On a finite
linear field $V=(V_i)$, the Gaussian divergence is
\[
 \delta V=\sum_i(G_iV_i-\partial_i V_i).
\]
Hence $\delta V_A=Q_{S(A)}$ for finite matrices, and
$\|\delta V_A\|_2^2=2\|S(A)\|_\HS^2$. It extends boundedly
to the completion of linear fields, with norm $\sqrt2$, and
agrees with the original closed divergence there. To check the last
assertion, for each smooth cylindrical test $f$, finite Gaussian
integration by parts gives
\[
 \E(f\,\delta V_{A^{(n)}})=\E\langle Df,V_{A^{(n)}}\rangle_H.
\]
Pass to the two $L^2$ limits. This is precisely the adjoint-domain
identity defining the original divergence, with a unique output.
Thus bounded completion is compatible with the original operator,
rather than defining a second unrelated divergence.

For symmetric $A$, the recovery field of $Q_A$ is $V_A$ and has
energy $\|V_A\|^2=\frac12\|Q_A\|_2^2$. On finite boxes this
is also $\frac12D Q_A$, since $D Q_A=2V_A$. The source's recovery
and symmetric-field unitary extend this identity to the full second
chaos. In this paper it suffices to retain the explicit coefficient
isometry and finite-box approximations. The matrix notation is a
human coordinate realization of those source constructions.

\section{Strong fourth-moment completion and bounded multiplication}
The finite upper bound applies to the symmetric difference of two
box matrices. Therefore
\[
 \|Q_{A^{(n)}}-Q_{A^{(m)}}\|_4
 \leq15^{1/4}\|Q_{A^{(n)}}-Q_{A^{(m)}}\|_2.
\]
The sequence is Cauchy in $L^4$. Its $L^4$ limit has the same
$L^2$ class as $Q_A$, because the probability measure gives
$\|X\|_2\leq\|X\|_4$. Completeness and uniqueness identify the
limits. Hence every $F\in\Hh_2$ belongs to $L^4$ and
\begin{equation}\label{eq:L4bound}
 \|F\|_4\leq15^{1/4}\|F\|_2.
\end{equation}
This also proves the upper bound in \eqref{eq:mainbounds} on the
full space. No continuity of squaring on unrestricted $L^2$ is used.

By H\"older and \eqref{eq:L4bound},
\[
 \|FG\|_2\leq\|F\|_4\|G\|_4
 \leq\sqrt{15}\|F\|_2\|G\|_2.
\]
The product is the actual almost-everywhere pointwise product, since
both factors have $L^4$ representatives. For box approximations,
\begin{align*}
 \|F_nG_n-FG\|_2
 &\leq\|F_n-F\|_4\|G_n\|_4
       +\|F\|_4\|G_n-G\|_4\longrightarrow0.
\end{align*}
In particular $F_n^2\to F^2$ strongly in $L^2$. Bilinearity on
equivalence classes is inherited from pointwise multiplication.

Take $F=G=G_0^2-1$. Its variance is $2$ and fourth moment is
$60$. The quotient $\|F^2\|_2/(\|F\|_2\|F\|_2)$ is
$\sqrt{60}/2=\sqrt{15}$. Thus the bilinear norm bound is exact,
and the upper fourth-moment coefficient is attained. The image
need not be quadratic: $(G_0^2-1)^2$ has a nonzero fourth-chaos
component, identified below.

\section{Degree projections and the completed cumulant}
For finite symmetric matrices the degree-two product component is
\begin{equation}\label{eq:mixedprojection}
 P_2(Q_AQ_B)=2Q_{AB+BA},\qquad P_2(Q_A^2)=4Q_{A^2}.
\end{equation}
One direct verification uses the degree-two product rule
\begin{align*}
 &(G_iG_j-\delta_{ij})(G_kG_l-\delta_{kl})\\
 &\quad=\text{degree-four Wick term}
 +\delta_{ik}(G_jG_l-\delta_{jl})
 +\delta_{il}(G_jG_k-\delta_{jk})\\
 &\qquad\quad+\delta_{jk}(G_iG_l-\delta_{il})
 +\delta_{jl}(G_iG_k-\delta_{ik})
 +\delta_{ik}\delta_{jl}+\delta_{il}\delta_{jk}.
\end{align*}
This identity follows by expanding the multivariate Hermite
generating function $\exp(\langle t,G\rangle-\|t\|^2/2)$ in its
two pairs of variables. Equivalently, subtract the listed lower
degree terms from the product and check its orthogonality to
degrees zero through three by Gaussian moments. Summing with
coefficients $A_{ij}B_{kl}$ gives \eqref{eq:mixedprojection}.
It gives $P_0(FG)=\langle F,G\rangle1$ as well.

For finite observables the product has only degrees $0,2,4$.
The sum of those three mutually orthogonal closed subspaces is
closed: if a sequence in their sum is Cauchy, each projected
component is Cauchy and converges in its closed subspace. Strong
product convergence therefore gives, for the full second chaos,
\begin{equation}\label{eq:decomposition}
 FG=P_0(FG)+P_2(FG)+P_4(FG).
\end{equation}
All other degree projections vanish. This is a decomposition in
the original Gaussian $L^2$, not a new multiplication internal
to $\Hh_2$.

For a finite symmetric $A$, \eqref{eq:mixedprojection} and
\eqref{eq:isometry} give
\[
 \|P_2(Q_A^2)\|_2^2=32\|A^2\|_\HS^2.
\]
Combining this with \eqref{eq:finitecumulant} proves
\eqref{eq:cumulant} on boxes. Strong $L^4$ approximation of $F$
and strong $L^2$ approximation of its square, followed by the
bounded orthogonal projection $P_2$, pass every term to the limit.
This proves \eqref{eq:cumulant} on the entire second chaos.
The argument uses finite-dimensional spectral theory only for the
finite fourth-moment calculation; no infinite spectral expansion
or trace-class hypothesis on the original quadratic matrix is needed.

Orthogonality in \eqref{eq:decomposition} also yields
\begin{equation}\label{eq:P4energy}
 \|P_4(F^2)\|_2^2
 =2\|F\|_2^4+\frac12\|P_2(F^2)\|_2^2.
\end{equation}
Indeed, $\|F^2\|_2^2=\|F\|_2^4+\|P_2(F^2)\|_2^2+
\|P_4(F^2)\|_2^2$, and \eqref{eq:cumulant} identifies the left
side. In particular,
\[
 \|P_2(F^2)\|_2^2\leq8\|F\|_2^4,\qquad
 2\|F\|_2^4\leq\|P_4(F^2)\|_2^2\leq6\|F\|_2^4.
\]
These are exact degree-energy consequences, not pointwise positivity
statements about the projected functions.

\section{Strictness without a uniform gap above three}
We prove that the second-degree square contraction of a nonzero
completed quadratic cannot vanish. This strictness does not follow
merely by taking limits of positive finite cumulants, since positive
numbers can tend to zero.

Represent $F=Q_A$ by its unique symmetric square-summable matrix.
For fixed $i,j$, its finite-box contraction satisfies, for $n>i,j$,
\[
 \frac14\bigl[\text{recovery of }P_2(Q_{A^{(n)}}^2)\bigr]_{ii}
 =\sum_{k<n}A_{ik}^2\geq A_{ij}^2.
\]
Here recovery sends $Q_C$ to $V_C$; an individual source coefficient
is a bounded linear functional on the completed field. Recovery
itself is bounded on the second chaos by the energy identity in
Section 3. Since $Q_{A^{(n)}}^2\to F^2$ in $L^2$, the coefficient
inequality passes to the limit. If $P_2(F^2)=0$, its recovered
diagonal coefficient is zero. Thus $A_{ij}^2=0$ for every fixed
pair $i,j$, so $A=0$ and $F=0$. This proves
\[
 F\ne0\ \Longrightarrow\ P_2(F^2)\ne0.
\]
Equation \eqref{eq:cumulant} now gives the strict lower inequality
in Theorem~\ref{thm:main}. Since every $F\in\Hh_2$ is centered,
no nonzero such variable has a Gaussian law: a centered Gaussian
of variance $\E F^2$ has fourth moment exactly $3(\E F^2)^2$.

Nevertheless there is no larger universal lower coefficient.
For $n\geq1$, set $T_n=\sum_{i=0}^{n-1}(G_i^2-1)$. The finite
moment calculation gives
\[
 \E T_n^2=2n,\qquad \E T_n^4=12n^2+48n,\qquad
 \frac{\E T_n^4}{(\E T_n^2)^2}=3+\frac{12}{n}\longrightarrow3.
\]
If a constant $c$ were a lower fourth-moment coefficient for every
quadratic in the completed space, these explicit tests would give
$c\leq3+12/n$ for every $n$, hence $c\leq3$. Thus strictness for
each nonzero target and optimality of the nonattained lower
coefficient coexist. No distributional limit theorem for $T_n$
is needed or asserted in this proof.

\section{Mixed square covariance by polarization}
We next prove the exact mixed identity \eqref{eq:polarized}.
Set $U=P_2(F^2)$, $V=P_2(G^2)$ and $W=P_2(FG)$.
Apply \eqref{eq:cumulant} to $F+G$ and $F-G$. Subtract twice
its versions for $F$ and $G$ from the sum of those two identities.
The fourth-moment contribution is $12\E(F^2G^2)$.
Writing $a=\|F\|_2^2$, $b=\|G\|_2^2$, $c=\langle F,G\rangle$,
the Gaussian subtraction contributes $-12ab-24c^2$.
On the contraction side,
\[
 P_2((F\pm G)^2)=U+V\pm2W.
\]
The parallelogram identity then gives the remaining contraction
contribution $6\langle U,V\rangle+12\|W\|_2^2$.
Dividing by $12$ proves \eqref{eq:polarized}.
All products and fourth moments exist by Section 4, so polarization
is valid on the completed space.

It remains to show $\langle U,V\rangle\geq0$.
For finite symmetric $A,B$, \eqref{eq:mixedprojection} gives
\begin{align}\label{eq:middlepair}
 \langle P_2(Q_A^2),P_2(Q_B^2)\rangle
 &=32\tr(A^2B^2)=32\|AB\|_\HS^2\geq0.
\end{align}
The matrix identity follows from
$\|AB\|_\HS^2=\tr((AB)^TAB)=\tr(BA^2B)=\tr(A^2B^2)$,
using symmetry and cyclicity of finite traces. Approximate $F,G$
by their finite boxes. Strong square convergence and boundedness
of $P_2$ make the two projected squares converge strongly, so
their inner products converge. Nonnegativity therefore passes
to arbitrary $F,G\in\Hh_2$. This proves the lower bounds and
the exact equality conditions in Theorem~\ref{thm:main}.

\section{Finite-matrix annihilation and the independence boundary}
For finite symmetric $A,B$, the preceding identities specialize to
\begin{align}\label{eq:matrixcov}
 \Cov(Q_A^2,Q_B^2)
 &=8(\tr AB)^2+16\|AB\|_\HS^2
       +8\|AB+BA\|_\HS^2.
\end{align}
Here $\langle Q_A,Q_B\rangle=2\tr AB$ and
$P_2(Q_AQ_B)=2Q_{AB+BA}$, while
\eqref{eq:middlepair} gives the middle term. Every term on the
right is nonnegative. It follows that
\begin{equation}\label{eq:finitecriterion}
 \Cov(Q_A^2,Q_B^2)=0\quad\Longleftrightarrow\quad AB=0.
\end{equation}
The reverse implication also uses $BA=(AB)^T=0$.
If $Q_A,Q_B$ are independent, their squares are independent and
their finite fourth moments imply factorization of the product of
squares. Equation \eqref{eq:finitecriterion} then forces $AB=0$.
This is the source-mapped necessary independence test.

For completeness, the classical converse is an elementary
finite-dimensional corollary, rather than an additional formal
declaration asserted in the source map. If $AB=0$, symmetry gives
$BA=0$ and $\operatorname{ran}A\perp\operatorname{ran}B$:
$\langle Ax,By\rangle=\langle x,ABy\rangle=0$.
Choose an orthonormal basis adapted to these two orthogonal ranges
and their joint orthogonal complement. A symmetric matrix
annihilates the orthogonal complement of its range; hence $Q_A$
depends only on the Gaussian coordinates in the first range, and
$Q_B$ only on those in the second. Those coordinate vectors are
independent, so the two observables are independent. Thus
\eqref{eq:finitecriterion} is also equivalent to independence in
this finite centered Gaussian setting. This is classical
quadratic-form theory, not a newly claimed criterion.

Orthogonality of the targets is weaker. For
\[
 A=\begin{pmatrix}1&0\\0&-1\end{pmatrix},\qquad
 B=\begin{pmatrix}0&1\\1&0\end{pmatrix},
\]
$\tr AB=0$, so $\langle Q_A,Q_B\rangle=0$.
But $AB\ne0$, $AB+BA=0$, and $\|AB\|_\HS^2=2$.
Equation \eqref{eq:matrixcov} gives
$\Cov(Q_A^2,Q_B^2)=32>0$, so they are dependent. The point is
the retained contraction, not covariance of the unsquared targets.
This explicit example is a finite matrix substitution in the
mapped identities, rather than a new standalone formal theorem.

\section{Scope and formal correspondence}
All sources are frozen at the verified commit
\begin{center}\texttt{\small bab688bb6140cd6316c80680095a914ea0d346db}.\end{center}
The theorem map records exact declarations, full headers, source
hashes and baseline links. The original full-build and complete
transitive-axiom receipt and its referenced evidence hashes are
checked separately. No research source was edited and no competing
full research build was run.

The principal sources concern original linear-field divergence
completion and recovery onto the second chaos, finite Gaussian
matrix moments, strong box approximation, actual completed products,
degree support, exact cumulants, strict contractions and mixed square
covariance. The symmetric square-summable matrix model above is a
coordinate reconstruction of their original source-field unitary.
The direct $L^4$ Cauchy proof reorganizes the finite sharp bound;
the formal square construction instead first identifies an actual
$L^2$ square through bounded approximation. Both retain the same
actual target, original measure and operator normalization.

The paper's general completed statements use the projection
conditions in \eqref{eq:polarized}. The matrix-annihilation
independence statement is presented for finite symmetric matrices.
No general infinite-dimensional independence classification,
Gaussian approximation rate or nonlinear process identification is
silently added. A Wick covariance power alone does not identify a
fractional process, generator or stochastic calculus.

Lancaster's primary indexed Theorem 1 gives the classical
trace-cumulant antecedent; its full theorem comparison remains
incomplete. Ogawa's primary introductory centered theorem and
finite orthogonal-range discussion supply the independence context.
The formalization value here is the complete extension to original
Gaussian $L^2$ targets and products with sharp constants, together
with explicit nonattainment and retained contraction terms. None
of these classical mechanisms is represented as a priority claim.

\begin{thebibliography}{9}
\bibitem{Lancaster} H. O. Lancaster,
\emph{Traces and cumulants of quadratic forms in normal variables},
Journal of the Royal Statistical Society, Series B \textbf{16}
(1954), 247--254.
\href{https://doi.org/10.1111/j.2517-6161.1954.tb00167.x}
{doi:10.1111/j.2517-6161.1954.tb00167.x}.
\bibitem{Ogawa} J. Ogawa,
\emph{On the independence of quadratic forms in a non-central normal system},
Osaka Mathematical Journal \textbf{2} (1950), 151--159.
\href{https://doi.org/10.18910/8824}{doi:10.18910/8824}.
\end{thebibliography}
\end{document}
